d'un morphisme projectif d'espaces analytiques complexes, et s'appliquera sans doute également en caractéristique quelconque, une fois résolu le problème des opérations «~puissances extérieures~» dans la catégorie dérivée (cf.\ Exp.\ XIV, n\textsuperscript{os}~1, 2, 3). C'est l'absence d'une étude de cette question qui constitue sans doute la lacune la plus sérieuse dans ce Séminaire, dans l'optique même où nous nous y sommes placés.

On remarquera l'absence, dans la table des matières du présent Séminaire, de deux des participants mentionnés sur la page de garde. L'un, J.P. Jouanolou, avait pris une part active à l'élaboration technique de la première partie du séminaire, mais avait été empêché de prendre part aux exposés oraux~; on lui doit notamment l'assertion d'indépendance linéaire contenue dans l'important énoncé VI 1.1. donnant la structure de l'anneau @@M0@@ des classes de Modules localement libres sur un fibré projectif (qui n'était prouvé auparavant que lorsqu'on supposait que le schéma de base admet un Module inversible ample). L'autre, J.P. Serre, avait fait deux exposés oraux, logiquement indépendants du reste du Séminaire, et qu'il a par suite préféré publier sous forme d'article séparé (cf.\ J.P. Serre, Groupes de Grothendieck des schémas en groupes réductifs déployés à paraître dans Publications Mathématiques, n\textsuperscript{o}~34).

Comme dans les autres parties du Séminaire de Géométrie Algébrique du Bois Marie, les sigles SGA 1 à SGA 6 renvoient aux différentes parties dudit Séminaire, le chiffre romain suivant ce sigle indiquant le n\textsuperscript{o} de l'exposé, et les chiffres arabes qui le suivent correspondant à la numérotation intérieure de l'exposé en question~; pour les références intérieures au présent Séminaire, on omet le sigle SGA 6, en utilisant des références style I 4.9 (signifiant~: exposé I, énoncé 4.9). Le sigle EGA réfère aux Éléments de Géométrie Algébrique de J. Dieudonné et A. Grothendieck.

